\chapter{Pecahan Rasional}\label{ch:b1:fractions}

\omterm{def:b1:fractions:field}{Pecahan rasional} adalah hasil bagi \omterm{def:b1:poly:def}{polinomial}. Teorema pusat
bab pendek ini --- yaitu penguraian pecahan
parsial\index{penguraian pecahan parsial} --- memecah sebarang
hasil bagi semacam itu menjadi jumlah bata dasar $\frac{c}{(X - a)^k}$.
Di luar minat aljabarnya, ia mesin baku untuk
mengintegralkan fungsi rasional (\cref{ch:b1:integration}) dan untuk
menjumlahkan deret tertentu (\cref{ch:b1:series}).

\section{Lapangan $K(X)$}

\begin{definition}\label{def:b1:fractions:field}
Sebuah \emph{pecahan rasional}\index{pecahan rasional} atas $K$ ($= \R$
atau $\C$) adalah hasil bagi $F = \frac{A}{B}$ dengan $A, B \in K[X]$, $B \neq
0$; dan dua hasil bagi $\frac AB$ serta $\frac{A'}{B'}$ diidentifikasikan bila
$AB' = A'B$. Setiap pecahan mempunyai \emph{bentuk paling sederhana} dengan
$\gcd(A, B) = 1$, yang tunggal sampai konstanta. Dengan operasi yang wajar,
\omterm{def:b1:logic:sets}{himpunan} $K(X)$ berisi pecahan rasional menjadi sebuah \omterm{def:b1:structures:field}{lapangan}.

Sebuah \emph{kutub}\index{kutub} $F$ (dalam bentuk paling sederhana) adalah akar
$B$; sedangkan \emph{orde} sebuah kutub adalah \omterm{def:b1:poly:derivative}{kegandaannya} sebagai akar $B$.
Adapun \emph{derajat} $F$ adalah $\deg F = \deg A - \deg B \in \Z \cup
\{-\infty\}$.
\end{definition}

\begin{example}[Membaca kutub, orde dan derajatnya]\label{ex:b1:fractions:vocab}
Misalkan $F = \dfrac{X^3 - X}{X^4 - 2X^3 + X^2}$. Faktorkan kedua lapisnya:
pembilangnya $X(X-1)(X+1)$, penyebutnya $X^2(X-1)^2$; lalu coret
faktor persekutuannya $X(X-1)$:
\[
F = \frac{X + 1}{X(X - 1)} \quad\text{(bentuk paling sederhana).}
\]
\omterm{def:b1:fractions:field}{Kutubnya}: $0$ dan $1$, keduanya \emph{sederhana} --- ordenya terbaca pada
penyebut yang sudah disederhanakan, sehingga akar ganda yang tampak pada penyebut
aslinya tidak relevan. Derajatnya: $\deg F = 1 - 2 = -1$,
yang terlihat secara asimtotik ($xF(x) \to 1$ ketika $x \to \infty$). Adapun
derajatnya berperilaku seperti derajat \omterm{def:b1:poly:def}{polinomial} ($\deg FG = \deg F +
\deg G$, $\deg(F + G) \leq \max$), yaitu kaidah pembukuan yang dipakai
terus-menerus pada perburuan koefisien di bawah: setiap argumen ``limit
$xF(x)$'' sesungguhnya cacahan derajat yang menyamar.
\end{example}

\begin{proposition}[Bagian bulat]\label{prop:b1:fractions:integerpart}
Setiap $F = \frac AB \in K(X)$ secara tunggal berbentuk $F = E + \frac RB$ dengan $E
\in K[X]$ (yaitu \emph{bagian bulat}, atau bagian \omterm{def:b1:poly:def}{polinomial}, dari $F$) dan
$\deg R < \deg B$. Di sini $E \neq 0$ jika dan hanya jika $\deg F \geq 0$.
\end{proposition}

\begin{proof}
Pembagian Euclid $A = BE + R$ (\cref{thm:b1:poly:division}), lalu dibagi
oleh $B$. Ketunggalannya: jika $E + \frac RB = E' + \frac{R'}{B}$, maka $(E -
E')B = R' - R$ dengan $\deg(R' - R) < \deg B$, yang memaksa $E = E'$ lalu $R
= R'$.
\end{proof}

\begin{example}[Sederhanakan dulu, baru bagi]\label{ex:b1:fractions:integerpart}
Carilah bagian bulat $F = \dfrac{X^3 + 1}{X^2 - 1}$. Dengan \omterm{def:b1:arith:divides}{membagi}
secara membabi buta: $X^3 + 1 = (X^2 - 1)X + (X + 1)$, jadi $F = X +
\frac{X + 1}{X^2 - 1}$. Tetapi pecahannya belum disederhanakan: $X^3 + 1
= (X + 1)(X^2 - X + 1)$ dan $X^2 - 1 = (X+1)(X-1)$ berbagi faktor
$X + 1$, sehingga
\[
F = \frac{X^2 - X + 1}{X - 1} = X + \frac{1}{X - 1} :
\]
bagian bulatnya tetap sama, yaitu $X$, tetapi bagian pecahannya runtuh menjadi
satu bata saja, dan ``\omterm{def:b1:fractions:field}{kutub}'' di $-1$ sesungguhnya tak pernah menjadi \omterm{def:b1:fractions:field}{kutub}.
Sederhanakanlah selalu sampai bentuk yang terkecil sebelum berburu \omterm{def:b1:fractions:field}{kutub}: karena \omterm{def:b1:fractions:field}{kutub}
$F$ adalah akar penyebut yang \emph{sudah disederhanakan}. (Adapun bagian
bulatnya tak peka terhadap penyederhanaan itu, sebagaimana dijamin ketunggalan pada
\cref{prop:b1:fractions:integerpart}.)
\end{example}

\section{Penguraian pecahan parsial atas $\C$}

\begin{theorem}[Penguraian atas $\C$]\label{thm:b1:fractions:complex}
Misalkan $F = \frac AB \in \C(X)$ dalam bentuk paling sederhana, dengan $B = c\,(X -
a_1)^{m_1} \cdots (X - a_r)^{m_r}$. Maka $F$ dapat ditulis, secara tunggal,
\[
F = E + \sum_{i=1}^{r} \sum_{k=1}^{m_i}
\frac{c_{i,k}}{(X - a_i)^k},
\qquad E \in \C[X],\ c_{i,k} \in \C ,
\]
dengan $E$ bagian bulat $F$.
\end{theorem}

\begin{proof}
Menurut \cref{prop:b1:fractions:integerpart} kita boleh mengandaikan $\deg A < \deg
B$ lalu membuktikan penguraian jumlahnya dengan $E = 0$.

\emph{Memilah \omterm{def:b1:fractions:field}{kutubnya}.} Tulis $B = (X - a_1)^{m_1} B_1$ dengan
$B_1(a_1) \neq 0$. \omterm{def:b1:poly:def}{Polinomial} $(X-a_1)^{m_1}$ dan $B_1$ \omterm{cor:b1:arith:bezout}{saling
prima} (karena tak berakar persekutuan), jadi menurut Bézout di $\C[X]$ (lihat catatan pada
\cref{ch:b1:poly}) ada $U_0, V_0$ dengan $U_0 (X-a_1)^{m_1} +
V_0 B_1 = 1$; lalu dengan mengalikannya dengan $A$ dan menulis $U = AU_0$, $V =
AV_0$, kemudian \omterm{def:b1:arith:divides}{membaginya} dengan $B$:
\[
\frac AB = \frac{V}{(X - a_1)^{m_1}} + \frac{U}{B_1} .
\]
Syarat derajatnya dapat dipaksakan: bagilah $V$ dengan
$(X-a_1)^{m_1}$, katakanlah $V = (X-a_1)^{m_1}Q + V_1$ dengan $\deg V_1 <
m_1$, lalu serap hasil baginya ke dalam suku kedua ($U_1 = U +
QB_1$):
\[
\frac AB = \frac{V_1}{(X - a_1)^{m_1}} + \frac{U_1}{B_1},
\qquad \deg V_1 < m_1 ;
\]
lalu membandingkan derajatnya ($\deg A < \deg B$ dan $\deg V_1 < m_1$) memaksa
$\deg U_1 < \deg B_1$ juga. Mengiterasikannya pada $\frac{U_1}{B_1}$,
\omterm{def:b1:fractions:field}{kutub} demi \omterm{def:b1:fractions:field}{kutub}, menyusutkan segalanya menjadi kasus \omterm{def:b1:fractions:field}{kutub} tunggal di bawah.

\emph{\omterm{def:b1:fractions:field}{Kutub} tunggal.} Untuk $\frac{A}{(X-a)^m}$ dengan $\deg A < m$: jabarkan $A$
dalam pangkat $(X - a)$, $A = \sum_{j=0}^{m-1} \alpha_j (X - a)^j$
(yaitu ekspansi Taylor sebuah \omterm{def:b1:poly:def}{polinomial}, seperti pada bukti
\cref{prop:b1:poly:multiplicity}); lalu \omterm{def:b1:arith:divides}{membaginya} memberikan tepat bata
$\frac{\alpha_j}{(X-a)^{m-j}}$. Secara konkret, untuk $\frac{X^2 +
1}{(X+2)^3}$: dengan mensubstitusikan $X = Y - 2$,
\[
X^2 + 1 = (Y - 2)^2 + 1 = Y^2 - 4Y + 5 ,
\qquad\text{sehingga}\qquad
\frac{X^2+1}{(X+2)^3} = \frac1{Y} - \frac4{Y^2} + \frac5{Y^3}
\]
dengan $Y = X + 2$: jadi ketiga batanya muncul lewat pembagian belaka atas
ekspansi yang tergeser --- yaitu jalur tercepat setiap kali ada satu
\omterm{def:b1:fractions:field}{kutub} berorde tinggi, dan jalur yang dianjurkan untuk
\cref{exo:b1:fractions:3}.

\emph{Ketunggalannya.} Andaikan dua penguraian berimpit; maka
selisihnya menjadi kesamaan $0 = \sum_{i,k} \frac{d_{i,k}}{(X -
a_i)^k}$. Kalikan seluruhnya dengan $(X - a_1)^{m_1}$: setiap sukunya
memperoleh faktor yang lenyap di $a_1$ \emph{kecuali} suku dengan
$i = 1$, $k = m_1$, yang koefisiennya menjadi $d_{1,m_1}$ ditambah
suku yang memikul sekurang-kurangnya satu faktor $(X - a_1)$. Lalu menghitung nilainya di
$a_1$ (yang sah: setelah perkalian itu, tak ada \omterm{def:b1:fractions:field}{kutub} yang tersisa di
$a_1$) memberikan $d_{1,m_1} = 0$. Dengan koefisien teratasnya lenyap,
ulangi dengan $(X - a_1)^{m_1 - 1}$, dan seterusnya turun sampai $k = 1$;
lalu berpindahlah ke \omterm{def:b1:fractions:field}{kutub} berikutnya. Semua $d_{i,k}$ lenyap: jadi
penguraiannya tunggal.
\end{proof}

\begin{method}[Menghitung koefisiennya]\label{met:b1:fractions:compute}
Dalam praktik, hindarilah Bézout; padukan saja:
\begin{enumerate}
  \item \emph{penutupan untuk pangkat tertingginya}: koefisien
        $\frac{1}{(X-a)^m}$ (dengan $m$ orde \omterm{def:b1:fractions:field}{kutub} $a$) adalah
        \[
        c_{a,m} = \Bigl[\,(X - a)^m F\,\Bigr]_{X = a};
        \]
        dan untuk \omterm{def:b1:fractions:field}{kutub} yang \emph{sederhana} pada $F = \frac AB$, nilainya
        $\frac{A(a)}{B'(a)}$;
  \item \emph{perhitungan nilai} pada titik yang nyaman dan \emph{limit}
        $xF(x)$ ketika $x \to \infty$ untuk mengumpulkan hubungan linear bagi
        koefisien yang tersisa;
  \item \emph{simetri} paritas atau konjugasi, bila ada, untuk memangkas
        pekerjaannya menjadi separuh.
\end{enumerate}
\end{method}

\begin{remark}[Jebakan yang lazim dengan pecahan parsial]\label{rem:b1:fractions:pitfalls}
\begin{enumerate}
  \item \emph{Melewatkan bagian bulatnya.} Penguraian batanya
        berlaku untuk pecahan berderajat $< 0$; jadi ketika $\deg A \geq
        \deg B$, bagilah dulu
        (\cref{prop:b1:fractions:integerpart}), atau perburuan
        koefisiennya akan melahirkan pertentangan.
  \item \emph{Penutupan di luar jangkauannya.} Mengalikan dengan
        $(X - a)^k$ lalu menghitung nilainya di $a$ menghasilkan koefisiennya
        hanya untuk $k = m$, yaitu orde \emph{penuh} \omterm{def:b1:fractions:field}{kutub} itu; sedangkan
        koefisien berorde lebih rendah menuntut hubungan lain (limit atau
        perhitungan nilai) --- lihat \cref{ex:b1:fractions:multiple}.
  \item \emph{Lupa menyederhanakan.} \omterm{def:b1:fractions:field}{Kutub} terbaca pada bentuk yang
        \emph{sudah disederhanakan}; jadi faktor persekutuan antara pembilang
        dan penyebut melahirkan \omterm{def:b1:fractions:field}{kutub} hantu
        (\cref{ex:b1:fractions:integerpart}).
  \item \emph{Bentuk bata yang salah atas $\R$.} Di atas kuadrat yang
        tak tereduksikan, pembilangnya \emph{afin} ($\alpha X +
        \beta$), bukan konstanta; jadi menulis $\frac{c}{X^2 + 1}$ saja
        menghilangkan penyelesaian --- bentuk yang benar didiktekan
        \cref{thm:b1:fractions:real}, dan tak pernah diimprovisasi.
\end{enumerate}
\end{remark}

\begin{proof}[Bukti rumus kutub sederhananya]
Di dekat \omterm{def:b1:fractions:field}{kutub} sederhana $a$: $B = (X - a) Q$ dengan $Q(a) \neq 0$, dan $B' =
Q + (X - a) Q'$, sehingga $B'(a) = Q(a)$. Jadi nilai penutupannya adalah
$\frac{A(a)}{Q(a)} = \frac{A(a)}{B'(a)}$.
\end{proof}

\begin{example}\label{ex:b1:fractions:simple}
Uraikan $F = \dfrac{1}{X(X-1)(X-2)}$. Ada tiga \omterm{def:b1:fractions:field}{kutub} sederhana; lakukan penutupan
pada masing-masingnya:
\[
c_0 = \frac{1}{(0-1)(0-2)} = \frac12,
\quad
c_1 = \frac{1}{1 \times (1 - 2)} = -1,
\quad
c_2 = \frac{1}{2 \times 1} = \frac12,
\]
jadi $F = \dfrac{1/2}{X} - \dfrac{1}{X-1} + \dfrac{1/2}{X-2}$.
\emph{Periksa di $X = 3$:} secara langsung, $F(3) = \frac{1}{3\cdot2\cdot1} =
\frac16$; sedangkan dari penguraiannya, $\frac{1/2}{3} - \frac{1}{2} +
\frac{1/2}{1} = \frac16 - \frac12 + \frac12 = \frac16$.
\end{example}

\begin{example}[Kutub ganda]\label{ex:b1:fractions:multiple}
Uraikan $F = \dfrac{X}{(X-1)^2 (X+1)}$. Bentuknya:
$\frac{a}{(X-1)^2} + \frac{b}{X-1} + \frac{c}{X+1}$.
\begin{itemize}
  \item Penutupan pada \omterm{def:b1:fractions:field}{kutub} gandanya: $a = \bigl[\frac{X}{X+1}\bigr]_{X=1}
        = \frac12$.
  \item Penutupan di $-1$: $c = \bigl[\frac{X}{(X-1)^2}\bigr]_{X=-1} =
        -\frac14$.
  \item Limit $xF(x)$ di $\infty$: $0 = b + c$, jadi $b = \frac14$.
\end{itemize}
\[
F = \frac{1/2}{(X-1)^2} + \frac{1/4}{X-1} - \frac{1/4}{X+1} .
\]
\emph{Periksa di $X = 0$:} $F(0) = 0$ dan $\frac12 - \frac14 - \frac14
= 0$.
\end{example}

\section{Penguraian atas $\R$}

\begin{theorem}[Penguraian atas $\R$]\label{thm:b1:fractions:real}
Misalkan $F \in \R(X)$ dalam bentuk paling sederhana dengan penyebut
\[
B = c \prod_i (X - a_i)^{m_i} \prod_j (X^2 + p_j X + q_j)^{n_j}
\qquad (p_j^2 - 4q_j < 0).
\]
Maka $F$ terurai secara tunggal sebagai bagian bulatnya ditambah suku
\[
\frac{c_{i,k}}{(X - a_i)^k} \quad (1 \leq k \leq m_i),
\qquad
\frac{\alpha_{j,l}\, X + \beta_{j,l}}{(X^2 + p_j X + q_j)^{l}}
\quad (1 \leq l \leq n_j),
\]
dengan koefisien yang real.
\end{theorem}

\begin{proof}
Uraikan atas $\C$ (\cref{thm:b1:fractions:complex}). Karena $F$
real, koefisien di atas \omterm{def:b1:fractions:field}{kutub} $\conj a$ (pada setiap ordenya) adalah
\omterm{def:b1:complex:field}{konjugat} koefisien di atas $a$ (terapkan konjugasi pada
penguraiannya lalu panggil ketunggalannya). Kelompokkan setiap pasangan \omterm{def:b1:complex:field}{konjugatnya}:
\[
\frac{c}{(X - z)^l} + \frac{\conj c}{(X - \conj z)^l}
= \frac{c\,(X - \conj z)^l + \conj c\,(X - z)^l}{\bigl(X^2 - 2\Re(z)X
+ \abs z^2\bigr)^{l}},
\]
yang pembilangnya menjadi \omterm{def:b1:complex:field}{konjugat} dirinya sendiri, sehingga real, berderajat $\leq l$;
lalu memisahkan kelipatan kuadrat realnya menurunkannya ke derajat
$\leq 1$ pada setiap aras $l$ (yaitu induksi turun yang kecil). Adapun \omterm{def:b1:fractions:field}{kutub} real
mempertahankan koefisien realnya (karena konjugasi menetapkannya). Ketunggalannya
menyusul dari ketunggalan atas $\C$.
\end{proof}

\begin{example}[Menyaksikan pemasangan konjugatnya]\label{ex:b1:fractions:pairing}
Mekanisme buktinya, pada kasus terkecil: atas $\C$, \omterm{def:b1:fractions:field}{kutub}
$\frac1{X^2+1}$ adalah $\pm\iu$, dengan koefisien penutupan
$\frac1{2\iu}$ di $\iu$ dan $\frac1{-2\iu}$ di $-\iu$ ---
yang saling \omterm{def:b1:complex:field}{konjugat}, sebagaimana diramalkan teoremanya:
\[
\frac{1}{X^2 + 1}
= \frac{1/(2\iu)}{X - \iu} - \frac{1/(2\iu)}{X + \iu} .
\]
Dengan menggabungkannya kembali atas penyebut persekutuannya:
\[
\frac{1}{2\iu}\cdot\frac{(X + \iu) - (X - \iu)}{X^2 + 1}
= \frac{1}{2\iu}\cdot\frac{2\iu}{X^2 + 1} = \frac{1}{X^2+1} :
\]
bagian imajinernya saling hapus dan bata realnya muncul kembali utuh.
Untuk integran yang real kita biasanya tak pernah meninggalkan $\R$ --- tetapi ketika
menghitung nilai \emph{jumlah} pada titik kompleks (seperti dilakukan soal akhir
pekan dengan \omterm{def:b1:complex:unity}{akar satuan}), bata kompleks menjadi mata uang yang
wajar, dan pemasangan inilah kurs tukar antara kedua
penguraiannya.
\end{example}

\begin{example}\label{ex:b1:fractions:real}
Uraikan $F = \dfrac{4}{(X^2+1)(X-1)^2}$ atas $\R$. Bentuknya:
$\frac{aX + b}{X^2 + 1} + \frac{c}{(X-1)^2} + \frac{d}{X - 1}$.
Penutupan pada \omterm{def:b1:fractions:field}{kutub} gandanya: $c = \bigl[\frac{4}{X^2+1}\bigr]_{X=1} =
2$. Penutupan pada \omterm{def:b1:fractions:field}{kutub} kompleks $\iu$ (dengan pembilang di atas $X^2 + 1$
dihitung lewat penguraian kompleksnya, atau langsung): kalikan dengan
$X^2 + 1$ lalu ambil $X = \iu$:
\[
a\iu + b = \frac{4}{(\iu - 1)^2} = \frac{4}{-2\iu} = 2\iu ,
\]
jadi $a = 2$, $b = 0$. Limit $xF(x)$ di tak hingga: $0 = a + d$, sehingga $d
= -2$. Jadi
\[
F = \frac{2X}{X^2+1} + \frac{2}{(X-1)^2} - \frac{2}{X-1} .
\]
\emph{Periksa di $X = 0$:} $F(0) = 4$ dan $0 + 2 + 2 = 4$.
\end{example}

\begin{example}[Penyebut kubik, dari awal sampai akhir]\label{ex:b1:fractions:cubic}
Uraikan $F = \dfrac{1}{X^3 + 1}$ atas $\R$. Faktorkan dulu: $X^3 +
1 = (X + 1)(X^2 - X + 1)$, dengan kuadratnya berdiskriminan $-3 <
0$. Bentuknya: $\frac{a}{X+1} + \frac{bX + c}{X^2 - X + 1}$. Penutupan
pada \omterm{def:b1:fractions:field}{kutub} sederhana $-1$: $a = \bigl[\frac1{X^2 - X +
1}\bigr]_{X=-1} = \frac13$. Limit $xF(x)$ di tak hingga: $0 = a +
b$, jadi $b = -\frac13$. Perhitungan nilai di $X = 0$: $1 = a + c$, jadi $c =
\frac23$. Sehingga
\[
\frac{1}{X^3+1} = \frac13\Bigl(\frac{1}{X+1}
+ \frac{-X + 2}{X^2 - X + 1}\Bigr) ,
\]
yang dibenarkan di $X = 1$: ruas kirinya $\frac12$, ruas kanannya
$\frac13\bigl(\frac12 + 1\bigr) = \frac12$. Perhatikan penghematannya:
tiga bilangan yang belum diketahui, tiga fakta linear yang murah (satu penutupan, satu
limit, satu perhitungan nilai), tanpa penjabaran apa pun --- yaitu alur kerja
\cref{met:b1:fractions:compute} dalam bentuk murninya.
\end{example}

\begin{remark}[Untuk apa semua ini]
Begitu terurai, sebuah fungsi rasional terintegralkan suku demi suku: bata
$\frac{1}{(x-a)^k}$ mempunyai antiturunan dasar, dan bata
$\frac{\alpha x + \beta}{(x^2 + px + q)^l}$ menyusut menjadi $\ln$ dan
$\arctan$ (\cref{ch:b1:integration}). Adapun jumlah teleskopik adalah penerapan
baku yang lain (\cref{exo:b1:fractions:8}).
\end{remark}

\begin{remark}[Di mana bab ini dipakai]\label{rem:b1:fractions:whereused}
Pecahan parsial di atas segalanya adalah langkah \emph{prapengolahan}: bab
pengintegralan (\cref{ch:b1:integration}) menyuapkan setiap integran rasional
lewat \cref{thm:b1:fractions:real} sebelum mengintegralkannya,
dan bab deret (\cref{ch:b1:series}) meneleskopkan suku rasional
persis seperti pada \cref{exo:b1:fractions:8} dan pada soal akhir
pekan di bawah --- yang mendorong tekniknya sampai ke
$\sum 1/k^2 = \pi^2/6$. Adapun turunan logaritmik $P'/P = \sum
m_i/(X - a_i)$ (\cref{exo:b1:poly:11}) muncul kembali setiap kali letak
akar ditelaah. Di luar jilid ini, penguraian
$1/\chi(X)$ untuk sebuah \omterm{def:b1:poly:def}{polinomial} karakteristik $\chi$ melandasi
perhitungan pangkat matriks dan transformasi Laplace pada jilid Tahun
ke-2: karena bata $\frac{c}{(X - a)^k}$ adalah bayangan aljabar
bagi penyelesaian $t^{k-1}\eu^{at}$ yang ditemui pada \cref{ch:b1:diffeq}.
\end{remark}

\begin{omfigure}
\begin{tikzpicture}[xscale=1.75, yscale=0.42]
  \draw[->, thin] (-2.7,0) -- (2.8,0) node[below] {$x$};
  \draw[->, thin] (0,-4.6) -- (0,4.9) node[left] {$y$};
  \draw[dashed, gray] (-1,-4.4) -- (-1,4.6);
  \draw[dashed, gray] (1,-4.4) -- (1,4.6);
  \draw[omThm, thin] (-2.6,1.5) -- (2.7,1.5)
    node[above right, font=\small, pos=0.02] {$y = \lambda$};
  \draw[omDef, thick, domain=-2.6:-1.31, samples=80, smooth]
    plot (\x, {1/(\x+1) + 1/\x + 1/(\x-1)});
  \draw[omDef, thick, domain=-0.83:-0.20, samples=120, smooth]
    plot (\x, {1/(\x+1) + 1/\x + 1/(\x-1)});
  \draw[omDef, thick, domain=0.20:0.83, samples=120, smooth]
    plot (\x, {1/(\x+1) + 1/\x + 1/(\x-1)});
  \draw[omDef, thick, domain=1.31:2.7, samples=80, smooth]
    plot (\x, {1/(\x+1) + 1/\x + 1/(\x-1)});
  \fill[omThm] (-0.715,1.5) circle (0.045 and 0.19);
  \fill[omThm] (0.40,1.5) circle (0.045 and 0.19);
  \fill[omThm] (2.3,1.5) circle (0.045 and 0.19);
\end{tikzpicture}

{\small Fungsi $F(x) = \frac1{x+1} + \frac1x +
\frac1{x-1}$ bertipe \cref{exo:b1:fractions:12}: ia turun tegas
pada setiap selang di antara \omterm{def:b1:fractions:field}{kutubnya} $-1, 0, 1$
(garis putus-putus). Setiap aras mendatar $\lambda > 0$ dipotong tepat
sekali per selang di sebelah kanan \omterm{def:b1:fractions:field}{kutub} pertamanya (titik yang ditandai):
jadi penyelesaian $F = \lambda$ berselang-seling dengan \omterm{def:b1:fractions:field}{kutubnya}.}
\end{omfigure}

\section{Latihan}

\begin{exercise}[$\star$]\label{exo:b1:fractions:1}
Uraikan atas $\R$: $\dfrac{1}{X^2 - 1}$;
$\;\dfrac{X}{X^2 - 3X + 2}$; $\;\dfrac{X^2 + 1}{X(X-1)}$ \emph{(perhatikan
bagian bulatnya)}.
\end{exercise}

\begin{exercise}[$\star$]\label{exo:b1:fractions:2}
Uraikan atas $\R$: $\dfrac{1}{X(X^2 + 1)}$ dan
$\dfrac{X^3}{X^2 + X + 1}$.
\end{exercise}

\begin{exercise}[$\star$]\label{exo:b1:fractions:3}
Uraikan $\dfrac{1}{X^2(X - 1)}$ dan $\dfrac{X + 1}{(X - 1)^3}$
\emph{(untuk yang kedua, substitusikan $Y = X - 1$)}.
\end{exercise}

\begin{exercise}[$\star\star$]\label{exo:b1:fractions:4}
Uraikan atas $\C$, lalu atas $\R$: $\dfrac{1}{X^4 - 1}$.
\end{exercise}

\begin{exercise}[$\star\star$]\label{exo:b1:fractions:5}
Uraikan atas $\R$: $\dfrac{X^2}{(X^2 + 1)^2}$, lalu simpulkan sebuah
antiturunan $x \mapsto \dfrac{x^2}{(x^2+1)^2}$ bila diberikan
$\int \frac{\dd x}{(x^2+1)^2} = \frac12\bigl(\arctan x +
\frac{x}{x^2+1}\bigr) + C$.
\end{exercise}

\begin{exercise}[$\star\star$]\label{exo:b1:fractions:6}
Untuk $n \in \N^*$, uraikan $F_n = \dfrac{n!}{X(X+1)\cdots(X+n)}$
\emph{(dengan \omterm{def:b1:fractions:field}{kutub} sederhana di $0, -1, \dots, -n$; pakai rumus penutupannya
lalu kenali \omterm{def:b1:counting:objects}{koefisien binomialnya})}.
\end{exercise}

\begin{exercise}[$\star\star$]\label{exo:b1:fractions:7}
Dengan memakai kesamaan pada \cref{exo:b1:poly:11} untuk $P = X^n - 1$,
buktikan bahwa
\[
\sum_{k=0}^{n-1} \frac{1}{X - \omega^k} = \frac{n X^{n-1}}{X^n - 1},
\qquad \omega = \eu^{2\iu\pi/n},
\]
lalu hitung nilai kedua ruasnya di $X = 2$ untuk $n = 4$ sebagai pemeriksaan.
\end{exercise}

\begin{exercise}[$\star\star$]\label{exo:b1:fractions:8}
Uraikan $\dfrac{1}{k(k+1)(k+2)}$ lalu hitunglah
\[
S_n = \sum_{k=1}^{n} \frac{1}{k(k+1)(k+2)},
\qquad\text{lalu}\qquad
\lim_{n \to \infty} S_n .
\]
\end{exercise}

\begin{exercise}[$\star\star\star$]\label{exo:b1:fractions:9}
Misalkan $P \in \R[X]$ \omterm{def:b1:poly:def}{monik} berderajat $n$ dengan $n$ akar real yang berbeda
$x_1 < \dots < x_n$. Buktikan bahwa
\[
\sum_{i=1}^{n} \frac{1}{P'(x_i)} = 0 \quad (n \geq 2),
\qquad
\sum_{i=1}^{n} \frac{x_i^{\,n-1}}{P'(x_i)} = 1 .
\]
\emph{Petunjuk: uraikan $\frac{X^m}{P}$ untuk $m \leq n - 1$ lalu perhatikan
peluruhan koefisiennya di tak hingga --- atau pakai \omterm{thm:b1:poly:lagrange}{interpolasi Lagrange}
(\cref{thm:b1:poly:lagrange}) atas monomial $X^m$ pada simpul
$x_i$.}
\end{exercise}

\begin{exercise}[$\star\star$]\label{exo:b1:fractions:10}
Uraikan $\dfrac{1}{X(X+1)^2}$ atas $\R$. Dengan menerima nilai
$\sum_{k \geq 1} \frac1{k^2} = \frac{\pi^2}6$ (yang dibuktikan pada soal akhir
pekan bab ini), simpulkan
\[
\sum_{n=1}^{\infty} \frac{1}{n(n+1)^2} = 2 - \frac{\pi^2}{6} .
\]
\end{exercise}

\begin{exercise}[$\star\star$]\label{exo:b1:fractions:11}
Uraikan atas $\R$: $\dfrac{1}{(X^2+1)(X^2+4)}$, lalu
$\dfrac{X^2}{(X^2+1)(X^2+4)}$. \emph{Petunjuk: kedua penyebutnya adalah
\omterm{def:b1:poly:def}{polinomial} dalam $X^2$: jadi uraikan $\frac1{(Y+1)(Y+4)}$ lebih dulu.}
\end{exercise}

\begin{exercise}[$\star\star\star$]\label{exo:b1:fractions:12}
(Persamaan sekuler) Misalkan $F = \sum_{i=1}^{r} \frac{c_i}{X - p_i}$
dengan $p_1 < p_2 < \dots < p_r$ real dan semua $c_i > 0$.
\begin{enumerate}
  \item Tunjukkan bahwa $F$ turun tegas pada setiap selang di
        daerah asalnya, lalu berikan limitnya di $\pm\infty$ dan pada kedua
        sisi setiap \omterm{def:b1:fractions:field}{kutubnya}.
  \item Simpulkan bahwa untuk setiap $\lambda > 0$, persamaan $F(x) =
        \lambda$ mempunyai tepat $r$ penyelesaian real, satu pada masing-masing
        selang $\intoo{p_i}{p_{i+1}}$ dan satu di luar $p_r$.
        (Persamaan semacam itu mengatur perturbasi nilai eigen; adapun
        teorema nilai antara di sini dipakai pada taraf sekolah menengah
        dan dibuktikan pada \cref{ch:b1:continuity}.)
\end{enumerate}
\end{exercise}

\section{Soal: Pecahan parsial sebagai sebuah mesin}

\begin{problem}\label{pb:b1:fractions:1}
Penguraian pecahan parsial tampak seperti pembukuan belaka; soal ini
menunjukkan bahwa ia sebuah mesin. Disuapi pecahan
$\frac1{X(X+1)\cdots(X+k)}$, ia meneleskopkan seluruh keluarga jumlah
menjadi bentuk tertutup; disuapi $\frac1{X^n - 1}$, ia menghasilkan
kesamaan trigonometri seperti
\[
\sum_{k=1}^{n-1}\frac{1}{\sin^2\frac{k\pi}{n}} = \frac{n^2-1}{3} ;
\]
dan, didorong satu langkah lagi, kesamaan itu memeras keluar salah satu
rumus paling termasyhur dalam matematika, yaitu rumus Euler
\[
\sum_{k=1}^{\infty}\frac1{k^2} = \frac{\pi^2}{6}
\]
--- yang di sini diperoleh tanpa apa pun di luar aljabar bab ini dan
trigonometri sekolah menengah.\index{teleskopik}\index{masalah Basel}
Di sepanjang soal ini, $\omega = \eu^{2\iu\pi/n}$; adapun limit barisan
dipakai pada taraf sekolah menengah (\cref{ch:b1:seq} meresmikannya).

\medskip
\noindent\textbf{Bagian I --- Teleskopnya.}
\begin{enumerate}
  \item Uraikan $\frac1{X(X+1)}$ lalu hitung $\sum_{n=1}^{N}
        \frac1{n(n+1)}$ secara eksak; lalu simpulkan bahwa jumlahnya menuju
        $1$.
  \item Hal yang sama untuk $\frac1{X(X+2)}$: tunjukkan $\sum_{n=1}^{N}
        \frac1{n(n+2)} = \frac12\bigl(\frac32 - \frac1{N+1} -
        \frac1{N+2}\bigr) \to \frac34$. (Dengan adanya celah, \emph{dua}
        suku batas bertahan pada masing-masing ujungnya.)
  \item Resmikan mekanismenya: jika $F(X) = G(X) - G(X+1)$ untuk
        suatu $G$ rasional yang tak berkutub di $\intco1{+\infty}$,
        maka $\sum_{n=1}^N F(n) = G(1) - G(N+1)$. Pulihkan
        nilai $\frac14$ pada \cref{exo:b1:fractions:8} dengan menunjukkan
        saksi $G$ untuk $F = \frac1{X(X+1)(X+2)}$.
  \item Buktikan teleskop faktorial yang umum: untuk $k \geq 1$,
        \[
        \frac{1}{X(X+1)\cdots(X+k)}
        = \frac1k\biggl(\frac{1}{X(X+1)\cdots(X+k-1)}
        - \frac{1}{(X+1)\cdots(X+k)}\biggr),
        \]
        lalu simpulkan
        \[
        \sum_{n=1}^{\infty}\frac{1}{n(n+1)\cdots(n+k)}
        = \frac{1}{k \cdot k!} .
        \]
        Periksa kasus $k = 2$ terhadap pertanyaan 3.
  \item Hitunglah nilai penguraian \cref{exo:b1:fractions:6}
        pada titik yang dipilih baik untuk membuktikan
        \[
        \sum_{j=0}^{n}(-1)^j\binom nj\,\frac1{j+1}
        = \frac1{n+1},
        \qquad
        \sum_{j=0}^{n}(-1)^j\binom nj\,\frac1{j+2}
        = \frac1{(n+1)(n+2)} .
        \]
\end{enumerate}

\medskip
\noindent\textbf{Bagian II --- Pecahan $1/(X^n - 1)$.}
\begin{enumerate}[resume]
  \item Tunjukkan lewat rumus penutupannya bahwa
        \[
        \frac{1}{X^n - 1}
        = \frac1n\sum_{k=0}^{n-1}\frac{\omega^k}{X - \omega^k} .
        \]
  \item Dua pemeriksaan kewajaran: periksa rumusnya secara langsung untuk $n =
        2$, lalu tunjukkan bahwa jumlah $n$ koefisiennya
        lenyap untuk $n \geq 2$ --- dan jelaskan mengapa ia harus demikian (tinjau
        $xF(x)$ ketika $x \to \infty$).
  \item Turunkan ulang lewat penutupan kesamaan pada
        \cref{exo:b1:fractions:7}:
        $\frac{nX^{n-1}}{X^n-1} = \sum_k \frac1{X - \omega^k}$.
  \item Kelompokkan \omterm{def:b1:fractions:field}{kutub} \omterm{def:b1:complex:field}{konjugatnya} untuk membuktikan penguraian realnya:
        dengan $\theta_k = \frac{2k\pi}n$,
        \[
        \frac{\omega^k}{X - \omega^k} +
        \frac{\omega^{n-k}}{X - \omega^{n-k}}
        = \frac{2\cos\theta_k\,X - 2}
        {X^2 - 2\cos\theta_k\,X + 1} ,
        \]
        lalu tuliskan penguraian real lengkap
        $\frac1{X^n-1}$ (bedakan $n$ ganjil dan $n$ genap).
  \item Khususkan pada $n = 4$ lalu cocokkan dengan
        \cref{exo:b1:fractions:4}.
\end{enumerate}

\medskip
\noindent\textbf{Bagian III --- Jumlah trigonometri, dan $\pi^2/6$
Euler.} Misalkan $P = 1 + X + \dots + X^{n-1}$, yang akarnya adalah
$\omega, \omega^2, \dots, \omega^{n-1}$ (semuanya sederhana).
\begin{enumerate}[resume]
  \item Dengan memakai $\frac{P'}P = \sum_{k=1}^{n-1}\frac1{X - \omega^k}$
        (\cref{exo:b1:poly:11}), buktikan
        \[
        \sum_{k=1}^{n-1}\frac{1}{1 - \omega^k} = \frac{n-1}2 .
        \]
  \item Buktikan $\dfrac1{1 - \eu^{\iu\theta}} = \dfrac12 +
        \dfrac\iu2\,\frac{\cos(\theta/2)}{\sin(\theta/2)}$ untuk
        $\theta \notin 2\pi\Z$ (lewat pemfaktoran setengah sudut,
        \cref{met:b1:complex:trig}), lalu simpulkan dari pertanyaan 11
        bahwa $\sum_{k=1}^{n-1}
        \frac{\cos(k\pi/n)}{\sin(k\pi/n)} = 0$ --- yang juga terlihat
        lewat simetri $k \leftrightarrow n - k$.
  \item Dengan menurunkan kesamaan pertanyaan 11 (yakni memakai
        $\bigl(\frac{P'}P\bigr)' = \frac{P''}P -
        \bigl(\frac{P'}P\bigr)^2$ yang dihitung di $X = 1$), buktikan
        \[
        \sum_{k=1}^{n-1}\frac{1}{(1 - \omega^k)^2}
        = \frac{(n-1)(5-n)}{12} .
        \]
  \item Dengan menulis $\cot t = \frac{\cos t}{\sin t}$, simpulkan dari
        pertanyaan 12--13 kedua bentuk tertutupnya
        \[
        \sum_{k=1}^{n-1}\cot^2\frac{k\pi}n = \frac{(n-1)(n-2)}3,
        \qquad
        \sum_{k=1}^{n-1}\frac1{\sin^2\frac{k\pi}n}
        = \frac{n^2-1}3 .
        \]
  \item Periksa kedua rumusnya dengan tangan untuk $n = 3$ dan $n = 4$.
  \item Buktikan ketaksamaan $\cot t < \frac1t < \frac1{\sin t}$
        untuk $t \in \intoo0{\frac\pi2}$ (dari $\sin t < t < \tan
        t$), lalu simpulkan, untuk $n = 2m + 1$ dan $1 \leq k \leq m$:
        \[
        \cot^2\frac{k\pi}n \;<\; \frac{n^2}{k^2\pi^2}
        \;<\; \frac1{\sin^2\frac{k\pi}n} .
        \]
  \item Jumlahkan ketaksamaan itu untuk $k = 1, \dots, m$ (dengan memakai
        simetri $k \leftrightarrow n - k$ untuk memaruhkan rumus
        pertanyaan 14) lalu apitlah:
        \[
        \sum_{k=1}^{\infty}\frac1{k^2} = \frac{\pi^2}6 .
        \]
\end{enumerate}

\medskip
\noindent\textbf{Bagian IV --- Bata berorde lebih tinggi.}
\begin{enumerate}[resume]
  \item Dengan mengkuadratkan penguraian $\frac1{X^2-1}$ lalu
        menguraikan ulang suku silangnya, buktikan
        \[
        \frac1{(X^2-1)^2}
        = \frac14\Bigl(\frac1{(X-1)^2} + \frac1{(X+1)^2}\Bigr)
        - \frac14\Bigl(\frac1{X-1} - \frac1{X+1}\Bigr),
        \]
        lalu periksa ia di $X = 0$.
  \item Padukan pertanyaan 18, teleskopnya, dan nilai Euler
        (pertanyaan 17) untuk membuktikan
        \[
        \sum_{n=2}^{\infty}\frac1{(n^2-1)^2}
        = \frac{\pi^2}{12} - \frac{11}{16} ,
        \]
        lalu benarkan nilainya secara numerik sampai tiga desimal.
  \item Turunkan kesamaan pertanyaan 8 untuk memperoleh sebuah
        bentuk tertutup bagi $\sum_{k=0}^{n-1}\frac1{(X -
        \omega^k)^2}$, lalu periksa ia di $X = 2$, $n = 2$.
  \item Buktikan bahwa untuk $k \geq 2$,
        \[
        \sum_{n=k}^{\infty}\frac{1}{\binom nk} = \frac{k}{k-1} .
        \]
        \emph{(Susutkan ke pertanyaan 4 dengan menulis $1/\binom nk$ memakai
        faktorial.)}
  \item Untuk $k = 3$, berikan jumlah parsial yang eksak
        $\sum_{n=3}^{N}\frac1{\binom n3}$ beserta limitnya.
\end{enumerate}

\medskip
\noindent\textbf{Bagian V --- Sintesis.}
\begin{enumerate}[resume]
  \item Sebagai perhitungan penutup, tuliskan penguraian real
        lengkap $\dfrac1{X^6 - 1}$ lalu periksa ia di $X =
        0$.
  \item Di mana persisnya soal ini memakai: (i) ketunggalan
        penguraiannya; (ii) \omterm{def:b1:complex:unity}{akar satuan} dari
        \cref{ch:b1:complex}; (iii) turunan logaritmik
        dari \cref{exo:b1:poly:11}? Satu kalimat untuk masing-masing.
  \item Sintesis, dalam satu paragraf pendek: satu kesamaan aljabar
        --- yaitu memecah pecahan menjadi bata --- melahirkan jumlah
        yang eksak, kesamaan trigonometri, dan $\pi^2/6$. Berilah komentar atas
        pembagian tugas antara aljabar (penguraian yang eksak dan
        berlaku di mana-mana) dan analisis (limit dan
        pengapitan), lalu tunjukkan tempat masing-masing benangnya
        diindustrialisasi: teleskop dan perbandingan pada
        \cref{ch:b1:series}, dan pengintegralan batanya pada
        \cref{ch:b1:integration}.

\end{enumerate}
\end{problem}
